\documentclass[11pt]{article}

\usepackage[T1]{fontenc}
\usepackage[utf8]{inputenc}
\usepackage{lmodern}
\usepackage[margin=1in]{geometry}
\usepackage{microtype}
\usepackage[numbers]{natbib}
\usepackage{hyperref}
\usepackage{enumitem}
\usepackage{titlesec}

\hypersetup{
  colorlinks=true,
  linkcolor=black,
  citecolor=blue,
  urlcolor=blue
}

\setlength{\parindent}{1.5em}
\setlength{\parskip}{0.25em}
\setlist[itemize]{leftmargin=1.5em,topsep=0.25em,itemsep=0.12em}
\titleformat{\section}{\normalfont\bfseries\large}{\thesection.}{0.6em}{}
\titleformat{\subsection}{\normalfont\bfseries\normalsize}{\thesubsection.}{0.6em}{}

\title{\textbf{Open Infrastructure for AI-Assisted Scientific Discovery}\\[0.4em]
\large Proposal for a \emph{Nature} Perspective}
\author{}
\date{April 2026}

\begin{document}
\maketitle

\section{Working title and target}
\emph{Open Infrastructure for AI-Assisted Scientific Discovery.} Target venue: a Perspective in \emph{Nature}. This memo is written for senior academic collaborators whom we are inviting to co-author or endorse the article, and should be readable in roughly ten minutes.

\section{Thesis}
AI systems are entering the cognitive core of research. Literature synthesis, hypothesis generation, experimental design, code, analysis, and writing are each being reshaped by capable models, and the workflows that combine these capabilities are coalescing into a new kind of scientific infrastructure. That infrastructure is overwhelmingly being built inside a small number of private platforms. If the default tooling, benchmarks, and interfaces are allowed to settle there, the norms that have made modern science legible---reproducibility, inspectability, contestability---cannot be assumed to survive the transition. A credible open stack, governed by scientific norms rather than commercial strategy, must exist while the design space is still fluid.

\section{Why a Perspective, why now}
Two observations motivate the timing. First, AI has begun to close, rather than merely advance, long-standing scientific problems---the AlphaFold result being the clearest instance to date \citep{jumper2021alphafold}---and broader surveys describe similar capability increases across representation learning, simulation, and experiment design \citep{wang2023discovery}. Second, these capability gains arrive against a well-documented background of declining idea productivity: sustaining a given rate of scientific progress now takes substantially more research effort than it did decades ago \citep{bloom2020ideas}. AI-assisted infrastructure could compress that gap, but the direction of compression depends on who holds the infrastructure. The Perspective format is well suited to this moment: it is forward-looking, argument-driven, and intended to coordinate a field rather than to report a single result.

\section{Proposed contribution: five pillars}
We propose that a credible public stack for AI-assisted science rests on five interlocking pillars, and that the Perspective is organized around them.

\begin{itemize}
    \item \textbf{Tacit-knowledge capture.} Much of what distinguishes skilled researchers---heuristics, failure modes, debugging instincts, practical judgment---rarely reaches published papers. A structured, versioned repository of this tacit layer is foundational.
    \item \textbf{Agent runtime for scientific workflows.} An execution environment with tool use, persistent memory, citation tracing, and long-horizon planning, designed for real research tasks rather than chat.
    \item \textbf{Scientific-ability benchmark.} Evaluation grounded in genuine research tasks: hypothesis generation, confounder detection, ablation design, methodological critique, and uncertainty calibration.
    \item \textbf{Scientific-reasoning model development.} Model improvement aimed at calibrated uncertainty, disciplined tool use, and truthful reasoning, rather than fluent persuasion.
    \item \textbf{Metascience and measurement.} The empirical study of how AI-assisted infrastructure actually reshapes scientific production---reviewer burden, citation networks, publication quality, replicability, and time-to-adoption. We treat this as a first-class pillar, not an afterthought: an open stack is only credible if its effects are measured, and the measurement is itself a public artifact.
\end{itemize}

\section{Empirical backing}
The article draws on existing literature---the productivity-slowdown results \citep{bloom2020ideas}, AlphaFold-class capability closures \citep{jumper2021alphafold}, and recent AI-for-science surveys \citep{wang2023discovery}---and on in-flight work by the proposing coalition. That in-flight work includes an open auto-research framework, a battery of paper-generation pipelines across scientific verticals, and a pair of social experiments designed as empirical probes of AI-mediated scientific contribution: one inside the existing peer-review pipeline (conferences, open review), and one on public software infrastructure, where agent-proposed and agent-reviewed features, issues, and pull requests provide naturalistic data on AI-mediated collaboration. Specific production counts and vertical-level results are reported in accompanying technical reports rather than in the Perspective itself.

\section{Coalition}
The initiative draws contributors from several leading research universities, with advisory and endorsement support from senior researchers in the agent, AI-for-science, and metascience communities. The article's author list is intended to reflect this coalition, with named co-authors for each of the five pillars and explicit acknowledgment of the broader contributor community.

\section{Timeline}
Assuming go-ahead in the coming weeks, we target a full draft within approximately four weeks of go-ahead and submission within approximately eight weeks. The article is timed to coincide with the public launch of the open platform and with several flagship conference submissions from the coalition that will serve as empirical anchors.

\section{On ``naive openness''}
We anticipate the natural critique that open infrastructure can be misused or poorly governed. The proposal is not for unconditional release. It is for inspectable, governed public infrastructure: infrastructure whose workflows, benchmarks, and models are open to scrutiny, and whose effects on scientific production are continuously measured and reported. Governance and metascience are load-bearing pillars of the design, not addenda---this is the point of treating measurement as co-equal with capability.

\bibliographystyle{unsrtnat}
\bibliography{refs}

\end{document}
